% Transcription — fonds Alexandre Grothendieck, shelfmark 37, batch 4 (pages 61-80).
% Established from the facsimile archives/batches/37/batch-04.pdf (a copy of
% 37.pdf fetched through the project's relay, grothendieck-relay.onrender.com;
% PDF page k = archivists' page 60+k, no cover sheet), per the skill
% transcribe-grothendieck. The folder is 99 pages in five batches (1-20, 21-40,
% 41-60, 61-80, 81-99); this is the fourth, done in parallel with the others.
%
% Pass: Opus 5.5 (claude-opus-5-5), 2026-09-26 — first pass, unchecked against the pages by a human.
%
% Pages skipped: 64 (a blank cream sheet). Page 72 is a tan cover inscribed
% in pencil, top right, « SGA 7 »: no \page marker, its words become a
% section heading with a \note (hand.md §5; the same treatment as the
% « SGA 1 », « SGA 4 », « SGA 5 », « SGA 6 » covers of batches 1 and 3).
% Pages noted only (not re-set): 65-71, 74-75.
% Pages transcribed: 61-63, 73, 76-80.
%
% Decisions, item by item, with the evidence:
% - Pages 61-63: his own manuscript notes on IHES letterhead, transcribed in
%   full. They continue the SGA 6 run opened by the cover « SGA 6 » (p. 55,
%   batch 3): p. 61 « Exp. X / Tenir compte conventions de Illusie (faire
%   relire par Berthelot) »; p. 62 « SGA 6 III 1, p. 15 », a footnote to add
%   (a flat morphism locally of finite presentation is perfect, a fortiori
%   pseudo-coherent); p. 63 a list of typescript pages each with an exposé
%   number (p. 43 (IX ...), p. 47 (IV), ... p. 75 (iv)), its object not said.
%   The section heading is therefore « SGA 6 (suite) », with a \note.
% - Pages 65-71: a typescript of seven leaves, « Commentaires Exp. X » (pp.
%   65-68 top) and « Commentaires Appendice » (pp. 68-71), a reader's
%   line-by-line remarks (« p.3, li.13 : ... ») on the typescript of an
%   Exposé X and of its appendix: missing references, typos, doubts (« ou
%   est-ce que je me canule ? »). WHICH SEMINAR: SGA 6, on the evidence of
%   the run it sits in (after the « SGA 6 » cover and p. 61, his note on
%   « Exp. X »), and of its internal references, cited by exposé number
%   alone, IV to VII and XIV (K.(X) « (IV 2.2 et 2.4) », « K.(X) naïf »,
%   « (VI 3.2) », « (V 3.10) », « Gr_top », gamma-filtration, Chern classes
%   c^i_chi), while SGA 5 IV and EGA are cited by name. WHOSE: DECISION — not
%   his; another person's typescript, one French \note per page, not re-set.
%   Evidence: the ink additions on pp. 65, 67, 69 and 71 (an insertion
%   « (V ...) », « ou Gamma(Z_X) », « N'y a-t-il pas une référence dans Bible
%   ou EGA 3 ? », « p. 13, li. 10 du bas : ... », « p. 16, li. 8 du bas : ...
%   ajouter "(cf. 7.16)" », « (7.15) », the note « (1) Est-ce que 7.2 ne se
%   généralise pas ? ») are in a small, regular, closed cursive plainly not
%   his fast, open hand of pp. 61-63 and 76-80; they read as the commenter
%   completing his own list, not as his reply to it. The reader is not
%   identified (p. 61 asks Berthelot to re-read Exposé X, which is
%   compatible with, not evidence for, Berthelot). No mark on the leaves is
%   attributable to him; the two pencil strokes in the left margin of p. 70
%   are given, their hand not established. The leaves are numbered by hand
%   « - 2 - » to « - 7 - » (the first unnumbered). The machine's missing
%   symbols (alpha, psi, chi, the arrows) are inked by the typist.
% - Page 73: a typed copy (thin paper, no signature) of a letter « Massy le
%   3.11.1970 » to « Cher Ferrand » asking whether Ferrand's « exposé SGA 6
%   XI » can be expected by the end of December, or else SGA 6 will be
%   published without it. DECISION: his letter, transcribed in full. Evidence:
%   the writer publishes SGA 6 and writes its introduction; same place, date
%   and form as the carbons of his letters of 3.11.1970 in batch 1 (pp. 3, 4,
%   15), one of which asks Deligne « As-tu eu des nouvelles de Ferrand pour
%   SGA 6 XI ? ». Letters clipped at the right edge are \supplied{}. It
%   concerns SGA 6 although it follows the « SGA 7 » cover.
% - Pages 74-75: a handwritten letter of P. Deligne (signed « P. Deligne »),
%   « Bures sur Yvette, le 21 octobre 1968 », blue ink on squared paper,
%   « Cher Grothendieck »: a proof, with the cotangent complex and GAGA, that
%   a compact complex analytic space X is a scheme iff X_red is. Another
%   person's letter: one \note per page, not re-set. His annotation, pencil,
%   top left of p. 74: « Joindre SGA 7 VI App. II ? » (first word doubtful),
%   matching his p. 76, « VI appendice : la lettre Deligne 2 pages ».
% - Pages 76-80: his own manuscript notes on IHES letterhead for SGA 7,
%   transcribed in full: p. 76 plan of appendices to VI, XXI, XXII (and a
%   possible XXIII); p. 77 (pencil) VI App.; p. 78 a cancelled item for VIII
%   and items for I and VI(?); pp. 79-80 lists of page/line places in IX and
%   XVIII, with bibliography items for IX.
%
% His pagination and numbered runs: none of his own on his leaves. His lists
% are keyed to the pages of typescripts not in this batch (p. 63: pp. 43-75;
% p. 79: pp. 1-169 of SGA 7 IX; p. 80: pp. 2-24 of SGA 7 XVIII). The
% commentary typescript is paginated by hand 2 to 7 (pp. 66-71), the first
% leaf unnumbered, and refers to pp. 1-37 of Exp. X and pp. 1-28 of the
% Appendice. Dates on the leaves: « le 21 octobre 1968 » (Deligne, p. 74);
% « 3.11.1970 » (p. 73).
%
% ALIGNMENT: batch 1 and batch 3 existed when this file was finished and
% were followed (cover treatment, \note format for other people's letters,
% his Massy letters of 3.11.1970 transcribed in full). Batches 2 and 5 did
% not exist.
%
% What the batch is about. The editorial preparation of SGA 6 and SGA 7:
% for SGA 6, Exposé X (conventions of Illusie, a re-reading by Berthelot, a
% reader's typed list of corrections to Exp. X and its appendix), a footnote
% for III 1 (flat + locally of finite presentation => perfect =>
% pseudo-coherent), and the letter chasing Ferrand's Exposé XI; for SGA 7,
% where Deligne's letter of 1968 (a compact analytic space is a scheme iff
% its reduction is) goes as an appendix to VI, appendices to XXI and XXII
% on crystalline cohomology (torsion), notation L Omega for the cotangent
% complex in VI App., the inertia group and its infinite p-Sylow (Serre),
% the « local invariant cycle problem » and Griffiths' conjecture, and lists
% of places to correct in IX and XVIII.
%
% Notation fixed once: roman exposé numbers as he writes them (he over- and
% underlines them; not reproduced); « l. » and « l.-n » for his line
% references (« -n » counting from the foot) kept as written.
%
% Hardest pages: 76 (the overwritten numeral of the first line, the dotted
% word « appendice »), 78 (a pencil-cancelled block and a crowded foot).
% Apparatus: 4 \ill{}, 17 \uncertain{} (counted on the finished file).
%
% Readings to check first: « AW » and the first sign on p. 63, and its last
% two numerals (« 74 », « iv »); the numeral before « 1.11 » on p. 76; « cot.
% rel. » and the index of L Omega on p. 77; the cancelled block and the Serre
% reference on p. 78; « 124 » and « livre » on p. 79; « Joindre » on p. 74.

\documentclass[11pt,a4paper]{article}
\input{../preamble/grothendieck}

\folder{37}
\batch{4}
\pages{61}{80}
\dating{1967-1971}
\watermark{Édition de démonstration}
\foldertitle{Rééditions SGA, EGA [SGA 4 à SGA 7, EGA I et EGA II] : tapuscrits et copies de tapuscrit annoté (s.d.), lettres (1967-1971, s.d.), notes manuscrites (s.d.)}

\begin{document}

\section*{SGA 6 (suite)}

\note{les pages 61 à 71 suivent le feuillet de garde « SGA 6 » (page 55, lot 3)
et les notes des pages 56 à 60 ; le titre est repris de ce feuillet}

\page{61}

\note{à l'encre, sur une feuille à en-tête de l'IHÉS ; rien d'autre sur la page}

\emph{Exp.} X \quad Tenir compte conventions de Illusie (faire relire par
Berthelot)

\page{62}

\note{à l'encre, sur une feuille à en-tête de l'IHÉS}

SGA \uncertain{6} III 1, p.~15, ajouter en footnote \struck{\ill{}} que un
morphisme plat loc.\ de prés.\ finie est parfait, et a fortiori
pseudo-cohérent.\note{le « 6 » est cerclé, et écrit d'un trait plus appuyé,
peut-être sur un autre chiffre ; « III » est surligné et souligné. Le petit
signe biffé est sous « SGA », au début de la deuxième ligne}

\page{63}

\note{à l'encre, sur une feuille à en-tête de l'IHÉS ; une liste de pages, chacune
suivie d'un numéro d'exposé entre parenthèses, les parenthèses laissées ouvertes
sur un blanc. Rien ne dit à quel tapuscrit elle renvoie}

\uncertain{$\alpha$} \uncertain{AW}\note{deux signes en tête, à gauche de la
liste : le premier ressemble à un $\alpha$ ou à une croix}

\begin{itemize}
\item p.~43 \quad (IX \ldots\ )
\item p.~47 \quad (IV \quad )
\item p.~51 \quad (IV \quad )
\item p.~52 \quad (IV \quad )
\item p.~56 \quad (IX \quad )
\item p.~59 \quad (IV \quad )
\item p.~70 \quad (IX \quad )
\item p.~\uncertain{74} \quad (IX \quad )\note{le chiffre romain est récrit sur
un premier tracé}
\item p.~75 \quad (\uncertain{iv} \quad )\note{en minuscules, à la différence des
autres}
\end{itemize}

\page{65}
\note{premier feuillet (non numéroté) d'un tapuscrit de sept feuillets intitulé
« Commentaires Exp.\ X », non reproduit : remarques d'un lecteur, page et ligne
par ligne, sur la frappe d'un exposé X, ici selon toute apparence celui de SGA 6
(voir la page 61 ; il cite par leur seul numéro les exposés IV à VII et XIV, à
part SGA 5 et EGA). Ce feuillet porte sur les pages 1 à 8 : références pour
$K_{\cdot}(X)$ (IV 2.2 et 2.4), « $K^{\cdot}(X)_{\text{naïf}}$ », (VI 3.2) si $X$
est quasi-compact, la numérotation du Cor.\ 1.3.3, $\alpha^{*}(x) \in
\mathrm{Filt}^{i}(Z)$, et des références dans EGA IV (3.1.4, 5.6.3, 5.7.1) pour
« universellement caténaire » et « Cohen-Macaulay ». Ajouts manuscrits, d'une
petite main régulière qui n'est pas la sienne : un renvoi « (V \ldots) » inséré
au-dessus de « référence pour » (ligne p.~5, li.~8) ; une accolade qui insère « une
autre lettre, Z par exemple, » dans la remarque sur le Cor.\ 1.3.3 ; et, en pied,
une ligne ajoutée à la fin de la remarque sur la p.~8, commençant par un mot
biffé, qui donne pour « l'assertion entre parenthèses » la référence (EGA IV
6.3.7). Rien de sa main}

\page{66}
\note{même tapuscrit, feuillet paginé « 2 » à la main, non reproduit : pages 9 à
25 de l'exposé (« biéquidimensionnel » et EGA $0_{\mathrm{IV}}$ 14.3.3, la
fonction de dimension définie par un complexe dualisant, « polynôme numérique »,
formules de changement de base, polynôme de Snapper, l'inégalité $\alpha_1 +
\cdots + \alpha_d \leq d$ dans (4.3.5), l'équivalence numérique et les
homomorphismes $\mathrm{Gr}_i(X) \to \mathbb{Z}$, $\mathrm{Gr}^i(X) \to
\mathbb{Z}$ définis par le produit d'intersection) ; il s'achève sur « Suggestion
aux §§ 2 et 4 : ajouter en remarque (ou exercice) que », suite page 67. Les
$\alpha$ et les flèches sont tracés à la main par la dactylographe. Aucune marque
manuscrite}

\page{67}
\note{feuillet « 3 », non reproduit : fin de la suggestion (2.3.1 et 4.3.1 pour
$F$ parfait résultent de R-R), puis pages 26 à 35 (« commutativement annelé »,
$H^0(X,\mathbb{Z})$ au lieu de $\Gamma(\mathbb{Z})$, $M =
T(\mathcal{O}^{n_1},\ldots,\mathcal{O}^{n_r})$ et la thèse de Giraud,
l'injection canonique $\mathrm{Der}(\mathrm{Gl}(n)) \hookrightarrow
\mathrm{Sl}(n)$, « voisinage étale » au lieu de « voisinage fpqc »). Ajouts de
la même main que page 65 : « ou $\Gamma(\mathbb{Z}_X)$ » inséré au-dessus de la
remarque sur la p.~29, et, à la suite de celle sur la p.~32 : « N'y a-t-il pas
une référence dans Bible ou EGA 3 ? ». Rien de sa main}

\page{68}
\note{feuillet « 4 », non reproduit : fin des commentaires de l'exposé (p.~37),
puis « Commentaires Appendice », pages 1 à 10 de l'appendice (preuve de 7.4, le
$K^{\cdot}$ de l'énoncé 7.1, « $X_t$ et $X_s$ réguliers $\Rightarrow$ $X$
régulier », les flèches du diagramme de 7.3, (7.10.2 bis), références pour
$\mathrm{Gr}^{\cdot}_{\mathrm{top}}$ et la compatibilité aux filtrations, un
diagramme analogue à (7.10.1 bis) avec $\mathrm{Gr}^{\cdot}(\ )_{\mathbb{Q}}$ pour
$g$ projectif d'intersection complète). Aucune marque manuscrite}

\page{69}
\note{feuillet « 5 », non reproduit : pages 11 à 17 de l'appendice ($K^{\cdot}(T)$
et non $K_{\cdot}(T)$, une référence attendue dans l'Exp.\ V, la spécialisation et
le changement de trait, Krull-Akizuki, les fibres géométriquement régulières d'un
$X$ non plat, et, en 7.13.6, deux formules reliant $\psi_i$, les classes $c^i$ et
les classes $c^i_{\chi}$, avec renvoi à (V 6.5.3 et 6.6.4)). Ajouts de la même main que page 65 : « (7.15) »
au bout de la remarque sur la p.~12, li.~6 ; une ligne sur la p.~13, li.~10 du bas
(« (IV 2.12) au lieu de (III 2.12) ») ; une ligne sur la p.~16, li.~8 du bas
(ajouter « (cf. 7.16) » après « on trouve »). Rien de sa main}

\page{70}
\note{feuillet « 6 », non reproduit : pages 17 à 25 de l'appendice (XIV 5.1 et
5.2, « linéairement » et « rationnellement », $\mathrm{Gr}^{\cdot} \xrightarrow{\sim}
\mathrm{Gr}^{\cdot}_{\mathrm{top}}$ modulo torsion (VII), SGA 5 IV, EGA IV 2.8,
$\dim(Z_s) = \dim(Z_t)$, EGA II 7.1.9, « essentiellement de type fini » et
« universellement japonais », « dommage qu'on ne donne pas de démonstration »).
Deux traits au crayon dans la marge de gauche, l'un court devant la remarque
« p.~24, li.~3 : supprimer "localement" », l'autre long devant « p.~24, dernière
ligne » ; main non établie. Rien d'autre}

\page{71}
\note{feuillet « 7 », dernier, non reproduit : pages 25 à 28 de l'appendice (un
« hint », une caractérisation d'un homomorphisme (7.17.4.1), « dans le
normalisé », et la référence Abhyankar sur la résolution des singularités des
surfaces arithmétiques qui manque à la bibliographie). Ajouts de la même main
que page 65 : une ligne tapée dans l'interligne, « (dans la situation présente,
i.e.\ S local régulier de pt fermé s) », appelée par une accolade ; un appel
« (1) » et, en pied, la note « (1) Est-ce que 7.2 ne se généralise pas ? ».
Rien de sa main}

\section*{SGA 7}

\note{inscrit au crayon en haut à droite d'une chemise de papier bistre (page
72), qui ne porte rien d'autre ; elle ne reçoit pas de numéro de page. La lettre
à Ferrand qui suit (page 73) concerne pourtant SGA 6}

\page{73}
\note{copie dactylographiée, non signée, d'une lettre de lui, sur papier pelure ;
le bord droit est rogné, les lettres perdues sont restituées}

Massy le 3.11.1970

Cher Ferrand,

J'aimerais savoir si je peux compter sur ton exposé SGA 6 XI dans un avenir
assez rapproché (diso\supplied{ns} d'ici fin décembre). Dans le cas contraire, je
pense qu'il serait préférable que je publie SGA 6 sans l'exposé XI. Il serait
quand même raisonnable, aprè\supplied{s} le travail que tu t'es tapé, que tu en
fasses au moins un article d'exposition, et j'aimerais savoir alors où tu penses
le publier, pour que je puisse y référer dans l'introduction à SGA 6.

Bien cordialement

\page{74}
\note{lettre manuscrite de P. Deligne, à l'encre bleue sur papier quadrillé,
« Bures sur Yvette, le 21 octobre 1968 », « Cher Grothendieck », feuillet 1 sur
2, non reproduite. Il montre qu'il n'est pas nécessaire de « faire de savants
recollages de complexes cotangents » pour prouver la proposition : \emph{soit $X$
un espace analytique complexe compact ; pour que $X$ soit un schéma, il faut et
il suffit que $X_{\mathrm{red}}$ le soit}. On se ramène au cas d'un Idéal
$\mathcal{N}$ de carré nul définissant $X_0$ ; avec le complexe cotangent $K$ de
$X_0$, donné localement dans la catégorie dérivée, il compare les suites exactes
$0 \to H^1(X_0, \mathcal{H}om(K,\mathcal{N})) \to E(X_0,\mathcal{N}) \to
H^0(X_0, \mathcal{E}xt^1(K,\mathcal{N}))$ algébrique et analytique ; par GAGA les
flèches verticales extrêmes sont des isomorphismes, et il reste à montrer que
$\mathrm{Im}(\beta) \subset \mathrm{Im}(\alpha)$. De sa main, au crayon, en haut à
gauche : « \uncertain{Joindre} SGA 7 VI App.\ II ? » (« VI » et « II » surlignés et
soulignés). Voir sa note de la page 76 : « VI appendice : la lettre Deligne 2
pages »}

\page{75}
\note{même lettre, feuillet 2, non reproduit : l'élément $e$ de $H^0(X_0,
\mathcal{E}xt^1(K,\mathcal{N}))$ est représenté par des extensions locales
$U'_i$ qui se recollent à isomorphisme près ; l'obstruction à choisir des
isomorphismes $\eta_{ij}$ satisfaisant la condition de cocycle est $d_2\eta$, où
$d_2 : H^0(X_0, \mathcal{E}xt^1(K,\mathcal{N})) \to H^2(X_0,
\mathcal{H}om(K,\mathcal{N}))$ ; la même construction marche dans le cas
analytique, et l'on conclut par GAGA. « Rmq : idem pour les espaces algébriques,
espaces rigides analytiques, schémas formels », puis une seconde « Rmq »
commencée et biffée. Signé « P. Deligne ». Aucune marque de sa main}

\page{76}

\note{à l'encre, sur une feuille à en-tête de l'IHÉS ; « SGA 7 », à gauche, est
entouré d'un tracé léger en forme de goutte. Une accolade réunit les deux
premières lignes, une autre les deux dernières}

SGA 7

\begin{itemize}
\item \uncertain{VI} 1.11 : 1.14 après 1.15--1.22 ? Refaire 1.14 ?\note{le chiffre
romain en tête est surchargé, récrit d'un trait épais sur un premier tracé ;
l'accolade qui le joint à la ligne suivante suggère VI}
\item VI \quad appendice : la lettre Deligne 2 pages\note{voir les pages 74 et 75}
\item XXI \quad appendice : relations avec torsion en cohomologie cristalline
\item XXII (?) \quad Commentaires sur les relations avec coh.\ cristalline.
\end{itemize}

\marginal{faire un exposé XXIII ?}\note{dans la marge de gauche, en regard de
l'accolade qui réunit les lignes XXI et XXII}

\page{77}

\note{au crayon, en haut d'une feuille blanche}

SGA 7 VI App

après 1.2 harmoniser notation \uncertain{$L\Omega_{X/k}$} pour compl.\
\uncertain{cot.} rel.\note{l'indice de $L\Omega$, après $X/$, est une boucle qui
peut être un $k$ ; « cot. » peut se lire « coh. »}

mettre bibliographie commune : l'exposé et l'App.\note{« mettre » est repassé
d'un trait plus appuyé}

\page{78}

\note{à l'encre, sur une feuille à en-tête de l'IHÉS. Le premier bloc, jusqu'au
trait horizontal, est biffé de grandes boucles au crayon}

\struck{SGA 7 VIII}

\struck{Faux schémas abéliens (ou \uncertain{motifs} mixtes effectifs de poids
1)}\note{« motifs » est récrit sur un premier mot}

\struck{Éléments de N.S. et accoupl.\ \ill{} symétriques.}

I \quad Le groupe d'inertie I a un \add{p-s-gpe de Sylow infini}\note{inséré
au-dessus de la ligne suivante}

cf Serre \uncertain{Corps} \ill{} \uncertain{Coh.\ galois.}, p.~144, lemme
4.\note{ligne soulignée d'un trait qui court sous toute la largeur}

\uncertain{VI} \quad Traiter aussi le « local invariant cycle problem » pour
\struck{des \ill{}} \add{$H^1$} ;
\uncertain{sign.}\ la conjecture \struck{de} Griffiths en dim.\ quelc.\note{le
chiffre romain est surchargé, noirci ; « $H^1$ » est écrit au-dessus des mots
biffés ; la ligne finale est soulignée d'un long trait}

\page{79}

\note{à l'encre, sur une feuille à en-tête de l'IHÉS : une liste de pages et de
lignes (« l.$-n$ » compte à partir du bas), à gauche ; une liste de la
bibliographie, à droite}

SGA 7 IX \quad (iii) \quad l.$-5$

\begin{itemize}
\item 1 \quad l.~5
\item 33 \quad l.~10
\item 35 \quad l.~5
\item 37 \quad l.~8
\item 69 \quad l.~$-3$
\item 70 \quad l.~2
\item 73 \quad l.~12
\item 87 \quad l.~6 \quad \uncertain{livre} de Gabriel-Demazure
\item \uncertain{124} \quad l.~$-8$\note{le chiffre du milieu peut se lire 7 ;
l'ordre croissant de la liste fait préférer 124}
\item 128 \quad l.~4, l.~$-9$
\item 134 \quad l.~1
\item 146 \quad l.~5
\item 150 \quad l.~$-3$
\item 157 \quad l.~5
\item 159 \quad l.~11, l.~$-7$\note{le 5 de 159 est récrit sur un autre chiffre}
\item 169 \quad l.~4 \quad deux fois
\end{itemize}

\emph{Bibl.} \quad \uncertain{ajouter} [22 bis] \quad [13] \quad [24]

\page{80}

\note{à l'encre, sur une feuille à en-tête de l'IHÉS ; rien d'autre sur la page}

SGA 7 XVIII

\begin{itemize}
\item p.~2 \quad 1.
\item p.~7 \quad l.~$-7$
\item p.~8 \quad l.~17
\item p.~24 \quad l.~8
\end{itemize}

\end{document}
